\section{字母诗歌}
\subsection{第9篇（大衛）：称颂傳講神的審判的作為-憐憫,公義  9:1-20}
\begin{itemize}
\itemsep-.45em  
\item Aleph \<'a> \footnote{https://cnbible.com} I will praise [thee] O LORD with all heart I will show forth all Your marvelous works. I will be glad and rejoice in I will sing praise to Your name O You most High 我要一心稱謝雅威，我要傳揚你一切奇妙的作為。我要因你歡喜快樂； 至高者啊，我要歌頌你的名！\footnote{https://bible.fhl.net}
\item Beth \<b> are turned When my enemies back they shall fall and perish at Your presence. For you have maintained my right and my cause you sat on the throne judging right. 我的仇敵轉身退去的時候，他們因你的面而跌倒滅亡。因你已經為我伸冤，為我辨屈；你坐在寶座上，按公義審判。
\item Gimel \<g> have rebuked the nations you have destroyed the wicked their name you have put out forever and ever 你曾斥責列邦，你曾滅絕惡人；你曾塗消他們的名，直到永永遠遠。  
\item Daleth \<d> \\
\item Hey \<h> O You enemy have come to a perpetual end destructions in perpetual and cities you have destroyed is perished their memorial with them  仇敵完全毀滅，永遠荒廢。你拆毀他們的城邑，連他們的名號也都滅沒。
\item Waw \<w> But the LORD forever shall endure he has prepared for judgment his throne 惟雅威坐著為王，直到永遠； 祂已經為審判擺設祂的寶座。
\item Waw \<w> And he shall judge the world in righteousness he shall minister judgment for the peoples with equity 祂要按公義審判世界，按正直判斷萬民。
\item Waw \<w> also will be The LORD a refuge for the oppressed a refuge in times of trouble  雅威要給受欺壓的人作高臺，在患難的時候作高臺。
\item Waw \<w> and will put their trust in they who know Your name for not do forsaken those who seek you LORD 雅威啊，認識你名的人要倚靠你，因你沒有離棄尋求你的人。
\item Zqyin  \<z> Sing praises to the LORD that dwells in Zion declare among the peoples his doings. When he makes inquisition for blood, he remember not do he forgets the cry afflicted. 應當歌頌居於錫安的雅威，將祂所行的傳揚在萬民中。那位追討血(債)的，祂記念他們(指受屈的人)，不忘記困苦人的哀求。
\item Het/Ceth  \<x> Have mercy O LORD consider my trouble from those who hate me you that lift me up from the gates of death. To the end that I may show forth all Your praise in the gates of the daughter of Zion I will rejoice in Your salvation 雅威啊，…求你憐恤我，(…處填入末行) 看那恨我的人所加給我的苦難，你是從死門把我提拔起來的；好叫我述說你一切(應得)的讚美；我必在錫安(原文是女子錫安)的城門因你的救恩歡樂。
\item Tet  \<T> are sunk down The nations in the pit [that] they made In the net that they hid taken is their own foot. is known The LORD [by] the judgment [which] he executes In the work of his own hands is snared the wicked Higgaion Selah. 他們(指外邦人)陷在自己所掘的坑中，他們的腳(落)在自己暗設的網羅裡。雅威已將自己顯明了，祂已施行審判；惡人被自己手所做的纏住了。(響聲，細拉)
\item Yodh \<y> shall be turned The wicked to Sheol all the nations that forget God. 惡人，…都必歸到陰間。(…處填入下行) 就是忘記上帝的外邦人，
\item Caph \<kk> For not always be forgotten For the needy the expectation humble [not] shall perish forever. Arise The LORD not do let prevail he be judged let the nations in Your sight. Put God them in hear themselves may know [that] the nations [to be but] men like Selah 窮乏人必不永久被忘，困苦人的指望必不永遠落空。雅威啊，求你起來，不容世人得勝！願列邦在你面前受審判！雅威啊，求你使他們恐懼，願外邦人知道自己不過是人。(細拉)
\end{itemize}

\subsubsection{詩歌如何教導今天我們的生活和禱告}
\begin{itemize}
\itemsep-.45em  
\item 1.	生命走入顺境，在得胜以后心中应该有清楚的认识，是神使我们得胜。我们旧造的人在取得成功后一般的反应是：（1）为自己取得的胜利而骄傲，因以为是自己努力的结果，到处炫耀自己的成就；（2）克服了难处，脱离了患难后，那个拦阻我们追求世界的虚荣或实现自我价值的东西被挪去了，于是我们将生命的轮子开足马力，去追求我们一直想要得到的目标；（3）或者是发现自己终于脱离了苦难，然后就尽情的在物质世界中享受各种所谓的快乐。
我们正确的做法是学习诗人在1-3节中所提的四件要做的事：将感谢，赞美归给神，在神面前欢喜快乐，并传扬神的作为。而能够做到这一点的关键操练就是学习诗人在困境中从依靠自己转向呼求依靠神，从而经历神的拯救。这样当我们脱离困境，进入顺境时，我们才会真心的将荣耀归给神。
\item 2.	在对待仇敌的态度上学习诗人，不用自己的手剪除仇敌，而是等候神的时间。
我们一般正常的做法是只要自己有能力，就用自己的手马上剪灭仇敌，或用自己的方法救自己脱离困境。
但诗人提到面对即使是自己有能力打倒的仇敌，正确的做法是：让仇敌因见神的面而绊跌灭亡。比如大卫的仇敌拿八羞辱他，大卫虽然有能力除掉他，但他听从了亚比该的劝诫没有杀拿八，从而看到神在拿八身上替自己伸冤；扫罗追杀他，但大卫两次都没有用自己的手拯救自己，而是依靠神。当神击打扫罗致死后，大卫因亲身经历到神信实的拯救，因而在便雅悯人士每咒骂他，约押不停的违抗他的命令杀了亚尼珥，亚马撒和押沙龙后，大卫依然不除灭他们而是继续信靠神。以至于所罗门学到了他在神面前等待的品格，没有马上杀死他们，而是等候神将他们的悖逆行为显现出来，然后才除灭这些恶人。
\item 3.	祷告方面的学习
\begin{itemize}
\itemsep-.45em  
\item a.	当神拯救诗人以后，他没有用笔墨述说因自己脱离仇敌和死亡的威胁而欢喜快乐，也不是因为仇敌被神消灭而快乐，反而是强调要因耶和华欢喜快乐(9:2)！
\item b.	大卫在这首诗篇中提到一个概念就是神救我们脱离一切苦难和困境的目的是要我们述说神一切的美德（9：14），而不是让我们继续在世界上实现自己的目标，或过世人所谓的快乐幸福的生活。并且大卫虽然早已不在了，但因为他信靠神，所有神今天依然透过他的口和笔，在向世人和我们宣扬神的在作为。
\end{itemize}
\end{itemize}

\subsubsection{诗篇格律-平行句}
\begin{itemize}
\item  同义平行 \\
（诗9:1）“我要一心称谢耶和华。我要传扬你一切奇妙的作为。”
（诗9:2）“我要因你欢喜快乐。至高者阿，我要歌颂你的名。”
  （诗9:19）“耶和华阿，求你起来，不容人得胜。愿外邦人在你面前受审判。”
（诗9:20）“耶和华阿，求你使外邦人恐惧。愿他们知道自己不过是人。细拉”
\item  延续平行\\
 （诗9:9）“耶和华又要给受欺压的人作高台，在患难的时候作高台。”
\end{itemize}

\subsubsection{诗篇特色}
\begin{itemize}
\item 结构方面 \\
：先赞美，后求告
（诗9:1-2）“（大卫的诗，交与伶长。调用慕拉便。）我要一心称谢耶和华。我要传扬你一切奇妙的作为。我要因你欢喜快乐。至高者阿，我要歌颂你的名。”
（诗9:19-20）“耶和华阿，求你起来，不容人得胜。愿外邦人在你面前受审判。耶和华阿，求你使外邦人恐惧。愿他们知道自己不过是人。（细拉）”

\item 修辞方面\\
：四个“我要”，三个“你曾”，四个“必”，两个“愿”。且这些相同的钥字，都出现在同一节，或者同一句话，或者同一段落内。
A、四个要：（诗9:1-2）“（大卫的诗，交与伶长。调用慕拉便。）我要一心称谢耶和华。我要传扬你一切奇妙的作为。我要因你欢喜快乐。至高者阿，我要歌颂你的名。”说明作者的心志。
B、三个“你曾”（诗9:5）“你曾斥责外邦。你曾灭绝恶人。你曾涂抹他们的名，直到永永远远。”说明上帝过去的作为。
C、四个“必”（诗9:14）“好叫我述说你一切的美德。我必在锡安城的门，因你的救恩欢乐。”（诗9:17-18）“恶人，就是忘记神的外邦人，都必归到阴间。穷乏人必不永久被忘，困苦人的指望，必不永远落空。”说明作者的确信。
D、两个“愿”（诗9:19-20）“耶和华阿，求你起来，不容人得胜。愿外邦人在你面前受审判。耶和华阿，求你使外邦人恐惧。愿他们知道自己不过是人。（细拉）”说明作者的祈愿。
\end{itemize}

\subsubsection{诗篇解析}
\begin{itemize}
\item 分段\\
诗篇第九篇一共二十节经文，（1-2）节为引言，（3-18）节为本论，（19-20）节为结论。

\item 解析\\
《赞美公义的审判者》（诗9篇）
一、立志赞美为引言（诗9：1-4）
1、立志赞美耶和华（诗9：1-2） 
（诗9:1-2）“（大卫的诗，交与伶长。调用慕拉便。）我要一心称谢耶和华。我要传扬你一切奇妙的作为。我要因你欢喜快乐。至高者阿，我要歌颂你的名。”（太六33）“你们要先求祂的国和祂的义，这些东西都要加给你们了。”这节金句正是今天这篇诗教导我们的真理，诗人在拯救没来临之先，就已献上称颂、感恩，并立志见证祂的救赎。诗人先想到神的荣耀，才想到自己的需要。赞美是表达我们对神的欣赏和对祂伟大的认识，也是对祂的神性的诸方面表示感谢。赞美也是把我们内在的态度变成外在的表白，藉此我们可以开扩对祂的认识。我们该作甚么才算是赞美神？赞美神的一种自然表露就是为祂作见证。当我们知道神的奇妙作为，很自然就想告诉别人，让他们也和我们一起赞美神。所以在（诗9：1-2），我们看到最多的是“称谢，传扬，喜乐，歌颂”。这些正是赞美当中的真意。

2、感谢神特殊干预（诗9：3-4）
（诗9：3-4）“我的仇敌转身退去的时候，他们一见你的面，就跌倒灭亡。因你已经为我伸冤，为我辨屈。你坐在宝座上，按公义审判。”神是我们的辩护者，替我们澄清批评、误会，在别人面前为我们辩护。我们可能面对许多不公义的事：我们可能被朋友误会，被仇敌诬告。我们可能对别人表示爱心，却没有被对方接纳。在工作和事奉上真正付上了代价，却没有适当的回报。我们的正确意见可能被忽视。但神是值得称赞的，因祂看见和记着我们所做的一切善行。应该由神来决定何时给我们何种报酬。假若我们不信祂会维护我们，我们就会忿恨，也很容易自怜；倘若我们信靠祂，就可以经历到由神而来的平安，不用再担心别人怎样理解和对待我们。

二、说明赞美的根据（诗9：5-6）
1、述说曾经的恩典（诗9：5）
2、指明仇敌的下场（诗9：6）
（诗9:5）“你曾斥责外邦。你曾灭绝恶人。你曾涂抹他们的名，直到永永远远。仇敌到了尽头。他们被毁坏，直到永远。你拆毁他们的城邑。连他们的名号，都归于无有。”诗篇9篇第5节，述说神曾经在祂选民身上的恩典。然后说呢，外邦的欺压者将受到严严的斥责，而以色列所有的仇敌将永远被遗忘。他们将要被埋没在那自夸为文明的废墟中。现在看似永恒的城市，将被连根拔起。像亚述，尼尼微大城。 诗人用预言的启示看到神为他的选民斥责恶人的时候。因信心可以认为那事已经成就了，所以就心中涌出赞美来。神选民的仇敌必被毁灭，就是他们的名号也无人纪念了。他们经耶和华审判的“风一吹便归无有，他的原处也不再认识他。”

三、神是主宰和保障（诗9：7-10）
1、神是公义审判的主（诗9：7-8）
2、神是慈爱怜悯的神（诗9：9-10）
（诗9:7-10）“惟耶和华坐着为王，直到永远。他已经为审判设摆他的宝座。他要按公义审判世界，按正直判断万民。耶和华又要给受欺压的人作高台，在患难的时候作高台。耶和华阿，认识你名的人要倚靠你。因你没有离弃寻求你的人。”对敌的将全都过去，惟耶和华坐着为王，直到永远，象祂一直以来那样正直可靠。祂在宝座发出灿烂的荣光，祂要以绝对的公义统治世界。人人都得到祂公平的对待。（徒一七30-31）保罗在雅典讲道的时候便使用第8节上半，解释将来负责审判的，是从死里复活的主耶稣基督：世人蒙昧无知的时候，神并不监察，如今却吩咐各处的人都要悔改。因为他已经定了日子，要借着他所设立的人，按公义审判天下，并且叫他从死里复活，给万人作可信的凭据。世上受欺压的民众要以祂为高台和可靠的避难所。所有认识祂的人都要倚靠祂，因为知道祂从来没有使祂的子民失望。神永远不会离弃那些信靠祂的人。神的应许，并非表示我们信靠祂就永不会经历痛苦和失败，而是表示神永远不离开我们。

四、呼吁赞美耶和华（诗9：11-12）
（诗9:11-12）“应当歌颂居锡安的耶和华，将他所行的传扬在众民中。因为那追讨流人血之罪的，他记念受屈的人，不忘记困苦人的哀求。”以色列不但要歌颂耶和华，还要完成向外邦众民传扬耶和华奇妙之拯救的使命，并且指出那追讨祂子民之血债的，不会不理他们的困苦，就是他们的祷告必蒙垂听。神不是单单住在锡安，而是时时刻刻住在每一处。不过以色列人敬拜的中心是耶路撒冷和漂亮的圣殿。神曾在会幕和所罗门所建造的圣殿中显现。犹太人把这位独一的真神从这个敬拜的中心地传向世界。

五、为困苦人的祈求（诗9：13-14）
（诗9:13-14）“耶和华阿，你是从死门把我提拔起来的。求你怜恤我。看那恨我的人所加给我的苦难。好叫我述说你一切的美德。我必在锡安城的门，因你的救恩欢乐。”遇到困境，人人都会祈求神来帮助，不同人祈求的原因也不尽相同。有些希望神帮助他们成功，令别人羡慕；另一些希望神帮助，使他们生活舒适。不过，大卫期望神帮助他使以色列恢复公义，同时可以向别人显明神的能力。当我们向神求助时，要反省一下自己的动机，是为了自己可以脱离痛苦和窘境，还是为了神的国度和尊荣？

六、讲述神拯救大能（诗9：15-18）
（诗9:15-18）“外邦人陷在自己所掘的坑中。他们的脚，在自己暗设的网罗里缠住了。耶和华已将自己显明了，他已施行审判。恶人被自己手所作的缠住了。（细拉）恶人，就是忘记神的外邦人，都必归到阴间。穷乏人必不永久被忘，困苦人的指望，必不永远落空。”跟着作者又向前跃到将来的时候，那时敌对闪族的外邦人，要掉进他们为犹太人所掘的坑中，并且踏进他们为神古时之民而设的网罗里。历史不断地重复——这就好象哈曼被挂在他为末底改而设的木架上一样。耶和华再次显明自己是施行公义的那位，祂要使恶人收取他们自己所种的恶果。神不能被嘲弄。大卫说恶人都必归到阴间指灵魂脱离肉体的状况，或指坟墓。根据上下文，这里应指地狱。这是忘记神的外邦人的终局。同样确定的是，穷乏人必不永久被忘。世人可能会漠视穷人的苦境，碾碎他们在世上的一线盼望，但神是软弱贫穷人的保护者，祂应许永不会忘记他们。充满罪恶的列邦忘记主，拒绝帮助祂的子民，将要被神审判。神知道我们的需要，也知道我们的灰心失望，（9：9、12）都说祂应许看顾我们，即使别人忘记了我们，祂仍然记念我们。

七、最后恳求神帮助（诗9：19-20）
（诗9:19-20）“耶和华阿，求你起来，不容人得胜。愿外邦人在你面前受审判。耶和华阿，求你使外邦人恐惧。愿他们知道自己不过是人。（细拉）”大卫对将来公义统治的思想，使他渴望这统治的临到。他祷告祈求耶和华起来阻止人的计划，和审判外邦人。他们站在全能的审判者面前的时候，会在恐惧中领会到自己不过是微不足道的人而已。我们要像大卫一样，赞美神公义的审判。阿们！
\end{itemize}


\subsection{第10篇：述说恶人用诡计毒害人，求神解救伸冤的祷告  10:1-18}
\begin{itemize}
\itemsep-.45em  
\item \underline{Lamed \<l>} Why O LORD stand you afar off hide yourself in times of trouble. the In pride wicked does persecute the poor let them be taken in the plots that they have imagined. For boasts the wicked on [the] desire of his heart and the covetous courses abhors the LORD. The wicked in the haughtiness of his countenance not do will seek not [after God] God all his thoughts。 prosper His ways at all times [are]  雅威啊，你為甚麼站在遠處？在患難的時候你為甚麼隱藏？惡人以驕橫追逼困苦人；願他們陷在自己所設的計謀裡。因為惡人以自己的心願自誇；他誇獎貪財的人，卻輕慢雅威。惡人鼻子高傲(，說)：祂(指雅威)必不追究；他一切所想的都以為沒有上帝。凡他所做的，時常穩固；
\item \underline{Mem \<mm>} for above Your judgments out of his sight [as for] all his enemies he puffs in them. He has said  to himself not do I shall be moved for generation and generation after not will be in adversity. of cursing his mouth is full and deceit and fraud Under his tongue [is] harm and vanity. He sits in the lurking of the villages. In the hiding does he murder 你的審判超過他的眼界。至於他一切的敵人，他都向他們噴氣。他心裡說：我必不動搖，世世代代不遭災難。他滿口是咒罵、詭詐、欺壓，他的舌底是毒害、奸惡。他在村莊埋伏等候，他在隱密處殺害
\item \underline{Nun \<nn>} the innocent. his eyes for the unfortunate are privately set. He lies in wait in a hiding as a lion in his lair, he lies in wait to catch 無辜的人,他的眼睛窺探無倚無靠的人；他埋伏在暗地，如獅子蹲在洞中。他埋伏，要擄去
\item \underline{Samech \<s>} 
\item \underline{Ain \<`>} the poor, he does catch the poor when he draws him into his net. He crouches humbles and may fall by his mighty that the poor. He has said to himself has forgotten God he hides. 困苦人；他拉網，就把困苦人擄去。他屈身蹲伏，無倚無靠的人就倒在他暴力之下。他心裡說：上帝竟忘記了,隱藏了；
\item \underline{Pe \<pp>} his face never see he will never 祂掩面永不觀看。
\item \underline{Tzaddi \<.s.s>} 
\item \underline{Koph \<q>} Arise O LORD O God lift up Your hand not do forget afflicted. Why how long scorn does the wicked. God, He has said to himself not do require. 雅威啊，求你興起！上帝啊，求你舉起你的手！不要忘記困苦人！惡人為何輕慢上帝，心裡說：你必不追究？
\item \underline{Resh  \<r>} have seen. for you harm and spite [it] behold to take [it] with Your hand himself to commits the poor of the fatherless You you are the helper. 其實你已經察看；你顧念憂患和愁苦，使它們在你管轄之中(原文是手中)。無倚無靠的人把自己交託給你；你向來是孤兒的幫助者。
\item \underline{Schin \<,s>} Break you the arm of the wicked and the evil [man] seek out his wickedness non [till] You find. The LORD [is] King forever and ever are perished the nations from His land. 願你打斷惡人的膀臂；至於壞人，願你追究他的惡，直到你找不著。雅威永永遠遠為王；列邦從祂的地已經滅絕了。
\item \underline{Tau \<t>} the desire of the humble you have heard LORD you will prepare their heart to hear you will cause Your ear. To vindicate the fatherless and the oppressed no more oppress he of the earth. 雅威啊，困苦人的心願，你早已聽見；你必堅固他們的心，也必側你的耳傾聽，為要給孤兒和受欺壓的人伸冤，使世上的人不再威嚇他們。
\end{itemize}
\subsubsection{教導}
困难是我们最不愿意遇到的事情，但又是我们无法回避的现实。这首诗歌教给我们一些关于如果渡过困境的智慧。
\begin{itemize}
\itemsep-.45em  
\item 在困境中我们可能会感觉不到神的帮助，或者离神很远，正如诗人在诗篇开头所述说的一样。这是每个神的儿女都会有的经历。
原来诗人也曾经有过类似的经历，所以我们不需要惊慌和害怕。我们感觉不到神同在的原因可能有很多：（1）我们因为高看自己的难处，以至于眼前的难处挡住了神的同在；（2）神是真的隐藏自己，为了锻造我们品格的益处；（3）我们的属灵生命不够成熟，辨别不出神一直同在的身影；（4）我们因为固执犯罪，而离开神的同在。诗人在这篇诗歌里，描述的是因为第二种原因感觉不到神的同在。所以如果我们确认自己是属于第二种情形，那么就要效法作者，相信神掌权，而带着盼望等候神。
\item 遇到恶人或恶事的时候既不是向恶人发作，也不是向第三者发泄心中的不满，而是应该向神控诉恶人的所为。
作者在18节的经文中，用了11节经文向神描述恶人的心思，言语和行为上的恶。从诗人的述说中可以看出这些人的所言所思所行实在是可恶，如果我们面临这样的恶人，可能会用“以其人之道治其人之身”的方法，直接向伤害我们的恶人反击回去，或者我们会向第三者述说恶人的所言所行，以寻求别人在行动或者道义上的声援和帮助，或者通过向别人抱怨的方法，以帮助化解自己心中的怒气。但诗人却选择了祷告的方法，详细向神描述恶人的所作所为。
\end{itemize}



\subsection{第25篇（大衛）： 诗人在患难中呼求神的怜悯和引导，并述说神的恩典 25:1-22}
\begin{itemize}
\itemsep-.45em  
\item \underline{Aleph \<'a>} unto you O LORD my soul do I lift up. O my God,\footnote{https://cnbible.com} 雅威啊，我的心仰望你。 我的上帝啊，\footnote{https://bible.fhl.net}
\item \underline{Beth \<b>} in you I trust let me not do be ashamed not do let triumph my enemies to. 我素來倚靠你；求你不要叫我羞愧，願我的仇敵不向我誇勝。
\item \underline{Gimel \<g>} Yes all that wait none on You be ashamed let them be ashamed that transgress cause outside. 凡等候你的必不羞愧，惟有那無故行奸詐的必要羞愧
\item \underline{Daleth \<d>} me Your ways O LORD, show me Your paths teach. 雅威啊，求祢將祢的道指示我，將祢的路指教我！
\item \underline{Hey \<h>} Lead me in Your truth. 求你以你的真理引導我，
\item \underline{Waw \<w>} and teach for you [are] the God of my salvation for on You do I wait all the day.指教我，因為你是救我的上帝。我終日等候你。
\item \underline{Zqyin  \<z>} remember Your tender mercies, O LORD, and Your covenant loyalty for they have been from of old they. 雅威啊，求你記念你的憐憫和慈愛，因為那些是亙古以來所常有的。
\item \underline{Het/Ceth  \<x>} the sins of my youth nor my transgressions not do Remember. According to Your covenant loyalty remember to you to the end that me for Your goodness O LORD. 求你不要記念我幼年的罪愆和我的過犯；雅威啊，求你因你的恩惠，按你的慈愛記念我
\item \underline{Tet  \<T>} Good and upright [is] the LORD, upon thus will he teach sinners in the way. 雅威是良善的，正直的，所以祂必教導罪人走正路。
\item \underline{Yodh \<y>} will he guide The meek in justice will he teach the meek his way. 祂要按公平引領謙卑人，將祂的道指教謙卑人。
\item \underline{Caph \<kk>} All the paths of the LORD [are] covenant loyalty and truth To those who keep his covenant and his testimonies. …雅威(對待他們)的方式是慈愛和信實。(…處填入下行)對遵守祂的約和祂法度的人，
\item \underline{Lamed \<l>} to the end that of For Your name, O LORD, and pardon my iniquity for [is] great it. 雅威啊，求你因你的名赦免我的罪，因為它(指我的罪)重大。
\item \underline{Mem \<mm>} What he who man [is] fears the LORD him shall he teach in the way [that] he shall choose. 誰敬畏雅威，祂(指雅威)必教導他當選擇的道路。
\item \underline{Nun \<nn>} His soul in prosperity shall dwell and his offspring shall inherit the earth.  他要安然居住，他的後裔必承受地土。
\item \underline{Samech \<s>} The secret of the LORD [is] with those who fear them his covenant him and he will show. 雅威的親密歸與祂的敬畏(指敬畏祂的人)， 要將自己的約指示他們。
\item \underline{Ain \<`>}  my eyes [are] ever toward the LORD, for he shall pluck my feet out out of the net my feet. 我的眼目時常仰望雅威，因為他必將我的腳從網裡拉出來。
\item \underline{Pe \<pp>} Turn you unto me and have mercy for [am] desolate and afflicted I [am]. 求你轉向我，憐恤我,因為我孤獨困苦。

\item \underline{Tzaddi \<.s.s>} The troubles of my heart are enlarged of my distresses [O] bring me out. 我心裡的愁苦甚多，求你帶領我脫離我的禍患。
\item \underline{Koph \<q>} 
\item \underline{Resh  \<r>} Look my affliction and my pain and forgive all my sins. Consider my enemies for they are many and they hate me with cruel hatred. 求你看顧我的困苦，我的艱難，求你赦免我一切的罪。求你察看我的仇敵，因為他們人多，並且以不公平的恨來恨我。
\item \underline{Schin \<,s>} O keep my soul and deliver let me not do be ashamed for I put my trust in. 求你保護我的性命，搭救我，使我不致羞愧，因為我投靠你。
\item \underline{Tau \<t>} Let integrity and uprightness preserve me for I wait on. Redeem, O God Israel, Out of all his troubles. 願純全、正直保守我，因為我等候你。上帝啊，求你救贖以色列脫離他一切的愁苦。
\end{itemize}

\subsubsection{教導}
\begin{itemize}
\itemsep-.45em  
\item 诗人在第3节提到凡等候神的必不羞愧。信仰初期，我们常常因怕祷告的事情超越神的掌管范围，从而神没有能力成就，犹豫不敢祷告，免得落入“羞愧”；或者因惧怕我们的信仰是假的，祷告就如对空说话而落入“羞愧”。这篇诗歌里，诗人教导我们，如何为我们的这种属灵盲点祷告，从而使我们的属灵生命成长。
另外一个导致我们祷告后，而可能落入“羞愧”的原因，是因为我们对神的认识不正确，因而做出错误的不符合神心意的祷告，比如“求神让我中大奖，我就奉献超过1/10的钱给教会”。但如果我们求神让我们不羞愧，神就会引导我们的心思意念，使得我们最终因为明白神的心意，而顺着圣灵的引导改变祷告方向，最终祷告的结果就是我们不会于落入羞愧。
\item 诗人在这篇诗歌里谈到“无故行奸诈的必要羞愧”。虽然我们肉眼看到的常常是，无故行奸诈的人计谋好像总是能得逞。但诗人在这里提醒我们，他们的结局是“必要羞愧”。不是靠我们征服他们，使他们羞愧；而是神会用各种方式，使得他们羞愧。（1）神将人的恶在人心中显明，使得人羞愧。比如神借着先知拿单用比喻责备大卫杀乌利亚并娶其妻子的恶行，使得大卫自感羞愧；（2）也可能是神在恶人身上施行审判，把恶人放在羞愧的地位上。如辱骂大卫的拿八，神击打他致死。（3）神借着另一个人所行的义，把人的恶显明出来，使做恶的人心中自觉羞愧。比如，神借着使得她玛怀孕生双子的方式，启示犹大的不守约，和践行诺言的她玛的义行。我们在这里要学习的功课就是：学会停止自己伸冤的行动，然后等待神的作为显明出来。
\end{itemize}


\subsection{第34篇（大衛）： 赞美神的救恩，要离恶行善靠主有福 34:1-22}
\begin{itemize}
\itemsep-.45em  
\item \underline{Aleph \<'a>} I will bless the LORD at all times, [shall] continually his praise [are] in my mouth \footnote{https://cnbible.com} 我要時時稱頌雅威,祂的讚美必常在我口中 \footnote{https://bible.fhl.net}
\item \underline{Beth \<b>} in the LORD shall make her boast My soul shall hear the humble [thereof] and be glad 我的心因雅威而感到驕傲，謙卑的人聽見就要喜樂
\item \underline{Gimel \<g>} O magnify the LORD with me and let us exalt his name together 你們要和我一同尊雅威為大，讓我們一同高舉祂的名。 
\item \underline{Daleth \<d>} I sought the LORD and he heard me from all my fears me And delivered 我曾尋求雅威，祂就應允我，救我脫離了我一切的恐懼。 
\item \underline{Hey \<h>} They looked unto and were lightened 他們仰望祂，就有光榮； 
\item \underline{Waw \<w>} and their faces were not do ashamed 他們的臉必不蒙羞。 
\item \underline{Zqyin  \<z>} This poor cried the LORD heard him out of all his troubles [him] and saved 這困苦人呼求，雅威便垂聽，救他脫離他一切的患難。
\item \underline{Het/Ceth  \<x>} encamps The angel of the LORD around those who fear him and delivers 雅威的使者在敬畏他的人四圍安營，要搭救他們。 
\item \underline{Tet  \<T>} O taste that see [is] good that the LORD blessed [is] the man [that] trusts in 你們要品嚐，就知道雅威是美善的；投靠祂的人有福了！
\item \underline{Yodh \<y>} O fear the LORD you his saints for not [there is] want to those who fear 雅威的聖民哪，你們當敬畏祂，因敬畏祂的一無所缺
\item \underline{Caph \<kk>} The young lions do lack but allow hunger they who seek the LORD not do want any good 少壯獅子還缺食忍餓，但尋求雅威的甚麼好處都不缺
\item \underline{Lamed \<l>} Come you children Listen to you the fear of the LORD I will teach 孩子們哪，你們當來聽我的話！我要將對雅威的敬畏教導你們 
\item \underline{Mem \<mm>} What man [is he that] desires life loves [many] days that he may see good 有何人喜好存活，愛慕長壽，得享美福？ 
\item \underline{Nun \<nn>} Keep Your tongue from evil and Your lips from speaking guile 你要保守你的舌頭不出惡言，你的嘴唇不說詭詐的話。 
\item \underline{Samech \<s>} Depart from evil and do good seek peace and pursue 要離惡行善，尋求和平，並追求它。 
\item \underline{Ain \<`>} The eyes of the LORD on [are] the righteous and his ears unto [are open] their cry 雅威的眼目向著(指看顧)義人，祂的耳朵向著(指聽)他們的呼求。 
\item \underline{Pe \<pp>} The face of the LORD [is] against those who do evil To cut of them from the earth the remembrance 雅威的臉敵對行惡的人，要從世上除滅他們的名號 
\item \underline{Tzaddi \<.s.s>} [The righteous] cry the LORD hears them out of all their troubles delivers 他們(指義人)呼求，雅威聽見了，便救他們脫離他們一切的患難。
\item \underline{Koph \<q>} [is] near The LORD to the broken hearted heart and such as be of a contrite spirit saves 雅威靠近心破碎的人，拯救靈性痛悔的人。 
\item \underline{Resh  \<r>} Many [are] the afflictions of the righteous him out of them all delivers but the LORD 義人的苦難很多，但雅威救他脫離這一切， 
\item \underline{Schin \<,s>} He keeps all his bones not one of them not do is broken 又保全他全身的骨頭，連一根也不折斷。 
\item \underline{Tau \<t>} shall slay the wicked Evil and they who hate the righteous shall be desolate. redeems The LORD the soul of his servants And none in him shall be desolate all of those who trust in 惡必害死惡人；恨惡義人的有罪。雅威救贖祂僕人的生命，凡投靠祂的，必不致定罪。 
\end{itemize}

\subsubsection{教導}
\begin{itemize}
\itemsep-.45em 
\item 诗人在这里不但自己依靠神，经历神的恩典，从而可以从心里发出真实的赞美，并且还呼吁带领众人与他一起依靠赞美神。
\item 诗人在这篇诗歌里谈到一个概念就是如果一个人“喜好存活，愛慕長壽，得享美福”，那么秘诀就是“禁止舌頭不出惡言，嘴唇不說詭詐的話”。把长寿和管住自己的嘴联系在一起，这和我们的观念大不一样。所以这里给我们在生活中的一个提醒就是，学习管住自己的舌头。就如雅各书所说：“若有人在話語上沒有過失，他就是完全人，也能勒住自己的全身”
\end{itemize}

\subsection{第37篇（大衛）： 不要因恶人而心懷不平而作惡，义人和恶人的结局 31:1-40}
\begin{itemize}
\itemsep-.45em  
\item \underline{Aleph \<'a>} neither do fret because of evildoers be not be You envious against the workers of iniquity. For like the grass they shall soon be cut down and as the green plant wither \footnote{https://cnbible.com}  不要因作惡的人怒火中燒，也不要妒忌那行不義的人。因為他們如草快被割下，又如綠色的青菜(快要)枯乾。\footnote{https://bible.fhl.net}
\item \underline{Beth \<b>} Trust in the LORD and do good [so] shall You dwell in the land and You shall be fed truly. and Delight in the LORD and he shall give to you the desires of Your heart. 你當倚靠雅威而行善，住在地上，以信實為糧；又要以雅威為樂，祂就將你心裡的請求賜給你。
\item \underline{Gimel \<g>} Commit to the LORD Your way and trust also in him and he [it] shall bring to pass. And he shall bring forth as the light Your righteousness and Your judgment as the noonday. 當將你的道路交託雅威，並倚靠祂，祂就必成全。祂要使你的公義如光發出， 使你的公平如正午。
\item \underline{Daleth \<d>} Rest in the LORD and wait patiently to fret not yourself ... because of him who prospers. in his way because of the man to pass wicked devices. 你當默然倚靠雅威，耐性(如待產般)等候他，不要因那道路通達的…怒火中燒。(…處填入下行)和那設惡謀的
\item \underline{Hey \<h>} Cease from anger and forsake wrath fret not yourself .. in any way to evildoing. for evildoers shall be cut off but those who wait on the LORD they shall inherit the earth. 當止住怒氣，離棄忿怒；不要怒火中燒，以致作惡。因為作惡的必被剪除；惟有等候雅威的，他們必承受地土。
\item \underline{Waw \<w>} For yet a little and not the wicked you shall diligently consider on for his place and it not. But the meek shall inherit the earth and shall delight themselves in the abundance of peace. 還有片時，惡人要歸於無有； 你就是細察他的住處，他也不存在。但謙卑的人必承受地土，以豐盛的平安為樂。
\item \underline{Zqyin  \<z>} plots The wicked against the righteous and gnashes on him with his teeth. The Lord shall laugh to for he sees that is coming his day. 惡人設謀害義人，又向他咬牙。主要笑他，因見他(受罰)的日子將要來到。
\item \underline{Het/Ceth  \<x>} the sword have drawn out The wicked and have bent their bow To cast the poor and needy To slay such as be of upright conversation. Their sword shall enter into their own heart and their bows shall be broken. 惡人已經刀出鞘，弓上弦，要打倒困苦窮乏的人，要殺害行為正直的人。他們的刀必刺入自己的心，他們的弓必被折斷。
\item \underline{Tet  \<T>} [is] man has better A little of the righteous than the riches wicked of many. For the arms of the wicked shall be broken and upholds the righteous but the LORD. 一個義人所有的雖少，強過許多惡人的富餘。因為惡人的膀臂必被折斷；但雅威扶持義人。
\item \underline{Yodh \<y>} knows The LORD the days of the upright and their inheritance forever shall be. not do be ashamed in the time in the evil and in the days of famine they shall be satisfied. 雅威知道完全人的日子，他們的產業要存到永遠。他們在患難的時候不致羞愧，在饑荒的日子必得飽足。
\item \underline{Caph \<kk>} for the wicked shall perish and the enemies of the LORD [shall be] as the fat of the pastures. they shall consume into smoke shall they consume away. 惡人卻要滅亡；雅威的仇敵要像草地的華美， 他們要消滅，在煙霧中消滅。
\item \underline{Lamed \<l>} borrows The wicked and not do pays but the righteous shows mercy and gives. for those blessed of him shall inherit the earth [they that be] and cursed of him shall be cut off. 惡人借貸而不償還；義人卻恩待人，並且施捨。蒙祂賜福的必承受地土；受祂詛咒的必被剪除。
\item \underline{Mem \<mm>} by the LORD The steps [good] of a man are ordered and in his way he delights. Though he fall not do be utterly cast down for the LORD upholds [him with] his hand. 一個人的腳步被雅威所立定，他的道路，祂也喜愛。他雖失腳也不至全身仆倒，因為雅威攙扶他的手。
\item \underline{Nun \<nn>} young I have been and [now] am old and not do seen the righteous forsaken nor his offspring begging bread. All [He is] ever merciful and lends and his offspring [is] blessed. 我從前年幼，現在年老，卻未見過義人被棄，他的後裔討飯。他終日恩待人，借給人；他的後裔也蒙福！
\item \underline{Samech \<s>} Depart from evil and do good and dwell forevermore. For the LORD loves judgment and not do forsakes his saints forever they are preserved but the offspring of the wicked shall be cut off. For the LORD loves judgment and not do forsakes his saints forever they are preserved but the offspring of the wicked shall be cut off. The righteous shall inherit the land and dwell in it forever. 你當離惡行善，就可永遠安居。因為，雅威喜愛公平；不撇棄祂的聖民，他們永蒙保佑，但惡人的後裔必被剪除。義人必承受地土，永遠定居
\item \underline{Ain \<`>} therein 在其上
\item \underline{Pe \<pp>} The mouth of the righteous speaks wisdom and his tongue talks of judgment. The law of his God [is] in his heart none shall slide of his steps. 義人的口談論智慧，他的舌頭講說公平。他的上帝的律法在他心裡，他的步伐總不滑跌。
\item \underline{Tzaddi \<.s.s>} watches The wicked upon the righteous and seeks to kill. The LORD not do leave him in his hand nor condemn him when he is judged.  惡人窺探義人，想要殺死他。雅威必不撇他在他(指惡人)的手中；當審判的時候，也不定他的罪。
\item \underline{Koph \<q>} Wait on the LORD and keep his way he shall exalt you to inherit the land are cut off when the wicked you shall see. 你當等候雅威，遵守祂的道，祂就抬舉你，使你承受地土；惡人被剪除的時候，你必看見。
\item \underline{Resh  \<r>} I have seen the wicked in great power and spreading himself bay tree like a green. And Yet he passed away see he not do I sought him but he could not do be found. 我見過惡人大有勢力，高聳如本地青翠的樹木。有人經過，不料，他不存在了；我尋找他，他卻無法被尋著。
\item \underline{Schin \<,s>} Mark 細察
\item \underline{Tau \<t>}  the perfect [man] and behold the upright for the end For the man [is] peace. But the transgressors shall be destroyed together the end of the wicked shall be cut off. But the salvation of the righteous [is] of the LORD [he is] their strength in time of trouble. And shall help the LORD them and deliver he shall deliver them from the wicked and save them. Because they trust in. 那完全人，觀看那正直人，因為(愛好)和平的人有(好)結局。至於罪人，必一同滅絕，惡人的未來必被剪除。但義人得救是出於雅威，祂在患難時作他們的避難所。雅威幫助他們，解救他們；祂解救他們脫離惡人，把他們救出來，因為他們投靠祂。
\end{itemize}












\subsection{第111篇：在众民的会中赞美耶和华的作为 111:1-10}
\begin{itemize}
\itemsep-.45em  
\item Praise you the LORD (HaliLu, Ya)\footnote{https://cnbible.com} 你們要讚美雅威（哈利路亚）\footnote{https://bible.fhl.net} 
\item \underline{Aleph \<'a>} I will praise the LORD with all [my] heart [in] and the congregation 我要一心稱謝雅威 
\item \underline{Beth \<b>} In the company of the upright 在正直人與會眾的大會中 
\item \underline{Gimel \<g>} Great [are] The works of the LORD 雅威的作為本為大   
\item \underline{Daleth \<d>} sought out by all those who have pleasure 為喜愛它們的人所考察 
\item \underline{Hey \<h>} [is] honorable and majestic His work 祂所行的是尊榮和威嚴 
\item \underline{Waw \<w>} and his righteousness endures forever 祂的公義存到永遠 \\
\item \underline{Zqyin  \<z>} to be remembered，He has made his wonderful works 祂行了奇事，使人記念 
\item \underline{Het/Ceth  \<x>} [is] gracious and full of compassion，the LORD 雅威有恩惠，有憐憫 
\item \underline{Tet  \<T>} food He has given to those who fear 祂賜糧食給敬畏祂的人 
\item \underline{Yodh \<y>} be mindful he will ever of his covenant 祂必永遠記念祂的約 
\item \underline{Caph \<kk>} the power of His works He has showed to His people 祂向祂的百姓顯出大能的作為 
\item \underline{Lamed \<l>} In giving to them the heritage of the nations 將列國賜給他們為業 
\item \underline{Mem \<mm>} The works of his hands [are] verity and judgment 祂手所行的是信實公平 
\item \underline{Nun \<nn>} sure All his commands 祂的訓詞都是確實的\\
\item \underline{Samech \<s>} They stand fast forever ever 是永永遠遠堅定的\\
\item \underline{Ain \<`>} [and are] done in truth and uprightness 是按信實正直設立的 
\item \underline{Pe \<pp>}  redemption He sent to His people 祂向祂的百姓施行救贖 
\item \underline{Tzaddi \<.s.s>} he has commanded forever his covenant 命定祂的約，直到永遠 
\item \underline{Koph \<q>} holy and reverend [is] his name 祂的名聖而可畏\\
\item \underline{Resh  \<r>}  [is] the beginning of wisdom The fear of the LORD 敬畏雅威是智慧的開端 
\item \underline{Schin \<,s>} a understanding good have all they who do 凡遵行它們(指祂的命令)的有美好的見識 
\item \underline{Tau \<t>} [his commandments] his praise endures forever 祂的讚美存到永遠 
\end{itemize}




\subsection{第112篇：敬畏遵守神命的福气 1-10}
\begin{itemize}
\itemsep-.45em  
\item Praise you the LORD. 你們要讚美雅威！
\item \underline{Aleph \<'a>} Blessed [is] the man [that] fears the LORD \footnote{https://cnbible.com}  敬畏雅威，…這人有福了！(…處填入下行)甚喜愛祂命令的，\footnote{https://bible.fhl.net}
\item \underline{Beth \<b>} the LORD in His commands [that] delights 甚喜愛祂命令的，
\item \underline{Gimel \<g>} mighty on earth shall be His offspring 他的後裔在世必強盛，
\item \underline{Daleth \<d>} the generation of the upright shall be blessed. 正直人的後代必蒙福。
\item \underline{Hey \<h>} Wealth and riches [shall be] in his house. 他家中有貨物，有錢財；
\item \underline{Waw \<w>} and his righteousness endures forever 他的仁義存到永遠。
\item \underline{Zqyin  \<z>} there rises in the darkness light for the upright 在黑暗中，有光向正直人發出；
\item \underline{Het/Ceth  \<x>} [he is] gracious and full of compassion and righteous 他有恩惠，有憐憫，有公義。
\item \underline{Tet  \<T>} A good man shows favor and lends 施恩與人、借貸與人、…的，這人事情順利；(…處填入下行)
\item \underline{Yodh \<y>} he will guide his affairs in judgment 秉公處理事務
\item \underline{Caph \<kk>} Surely For he will never not do be moved. 他不動搖，直到永遠； 
\item \underline{Lamed \<l>} remembrance in everlasting shall be the righteous  義人被記念，直到永遠。
\item \underline{Mem \<mm>} news of evil not do be afraid 他必不怕凶惡的信息；
\item \underline{Nun \<nn>} is fixed his heart trusting in the LORD 他的心堅定，倚靠雅威。
\item \underline{Samech \<s>} [is] established His heart not do be afraid 他的心確定，總不懼怕，
\item \underline{Ain \<`>} until that he see [his desire] on his enemies. 直到他看見他的敵人(遭報)。
\item \underline{Pe \<pp>} He has dispersed. he has given to the poor 他施捨，賙濟貧窮；
\item \underline{Tzaddi \<.s.s>} his righteousness endures forever 他的仁義存到永遠。
\item \underline{Koph \<q>} his horn shall be exalted in honor. 他的角必在榮耀中被高舉。
\item \underline{Resh  \<r>} The wicked shall see [it] and be grieved. 惡人看見便惱恨，
\item \underline{Schin \<,s>} with his teeth he shall gnash and melt away. 必咬牙而被融化；
\item \underline{Tau \<t>} the desire of the wicked shall perish. 惡人的心願要歸幻滅。
\end{itemize}

\subsection{第119篇： 神的法度/律例/话/典章 176}













\subsection{第143篇（大衛）： 靈裡枯竭的時侯，仰望神的慈愛 143:1-12}
\begin{itemize}
\itemsep-.45em  
\item  O LORD, Hear my prayer, give ear 雅威啊，求你聽我的禱告，側耳聽
\item \underline{Aleph \<'a>}\footnote{https://cnbible.com} unto my supplications 我的懇求， \footnote{https://bible.fhl.net}
\item \underline{Beth \<b>} in Your faithfulness answer, in Your righteousness. And not do enter into judgment with Your servant for not do be justified. For in Your sight all living. For has persecuted the enemy my soul. 憑你的信實和公義應允我。求你不要跟你的僕人進入審判關係，因為在你面前，凡活著的人沒有一個是義的。原來仇敵逼迫我，
\item \underline{Gimel \<g>} 
\item \underline{Daleth \<d>} he has struck to the ground my life
\item \underline{Hey \<h>} he has made me to dwell in dark dead as those who have been long.
\item \underline{Waw \<w>} and overwhelmed inside Therefore is my spirit inside me is desolate me My heart.
\item \underline{Zqyin  \<z>} I remember the days of old. I meditate on all Your works on the work of Your hands I muse. I stretch forth 
\item \underline{Het/Ceth  \<x>} 
\item \underline{Tet  \<T>} 
\item \underline{Yodh \<y>} my hands unto You. my soul 
\item \underline{Caph \<kk>} land as a thirsty
\item \underline{Lamed \<l>} to Selah 
\item \underline{Mem \<mm>} me speedily 
\item \underline{Nun \<nn>} 
\item \underline{Samech \<s>} 
\item \underline{Ain \<`>} Hear O LORD, fails my spirit not do hide 
\item \underline{Pe \<pp>} Your face from me lest I be like to those who go down to the pit. Cause me to hear in the morning Your covenant loyalty for in For I trust cause me to know the way wherein I should walk for unto I lift up my soul.
\item \underline{Tzaddi \<.s.s>}
\item \underline{Koph \<q>} 
\item \underline{Resh  \<r>} 
\item \underline{Schin \<,s>} 
\item \underline{Tau \<t>} 
\end{itemize}

\subsection{第145篇（大衛）： 凡有血氣的都稱頌傳講神的国和神手所做 1-21}
\begin{itemize}
\itemsep-.45em  
\item \underline{Aleph \<'a>} I will extol God, O King, and I will bless Your name forever and ever. \footnote{https://cnbible.com} 王, 我的上帝啊，我要尊崇你！我要永永遠遠稱頌你的名！ \footnote{https://bible.fhl.net}
\item \underline{Beth \<b>} Every day will I bless you and I will praise Your name forever and ever. 我要天天稱頌你，也要永永遠遠讚美你的名！
\item \underline{Gimel \<g>} Great [is] the LORD and to be praised greatly and his greatness not searchable. 雅威本為大，該受大讚美；其大無法測度。
\item \underline{Daleth \<d>} form generation to generation shall praise Your works and Your mighty acts shall declare. 這代要對那代頌讚你的作為，他們要傳揚你的大能。 
\item \underline{Hey \<h>} honor On the majestic of Your majesty and works of Your wondrous I will speak. 你威嚴榮耀的尊榮,和你奇妙的作為，我要默念。
\item \underline{Waw \<w>} and of the might of Your terrible acts [men] And shall speak Your greatness. I will declare. 你可畏之事的能力，人要傳講；你的偉大，我也要傳揚。
\item \underline{Zqyin  \<z>} the memory of Your great goodness. They shall abundantly utter and of Your righteousness shall sing. 他們要將你可記念的大恩傳開來，並要歌唱你的公義。
\item \underline{Het/Ceth  \<x>} [is] gracious and full of compassion. The LORD slow to anger and of great covenant loyalty. 雅威有恩典，有憐憫，不輕易發怒，大有慈愛。
\item \underline{Tet  \<T>} [is] good The LORD to all and his tender mercies over all [are] His works. 雅威善待萬有，祂的慈悲覆庇祂一切所造的。
\item \underline{Yodh \<y>} shall praise you O LORD, All Your works and Your saints shall bless. 雅威啊，你一切所造的都要稱謝你，你的聖民也要稱頌你；
\item \underline{Caph \<kk>} of the glory of Your Kingdom, They shall speak and of Your power talk. 他們要傳講你國度的榮耀，談論你的大能，
\item \underline{Lamed \<l>} to make atonement to the sons of men, his mighty acts and the majestic majesty of Your Kingdom. 好叫世人知道祂的大能，並祂國度威嚴的榮耀。
\item \underline{Mem \<mm>} Your kingdom kingdom, [endures] throughout all [is] an everlasting and Your dominion throughout all [endures] throughout generation and generation. 你的國是永遠的國！你的權柄存到萬代！
\item \underline{Nun \<nn>} 
\item \underline{Samech \<s>} upholds The LORD all that fall and raises up all [those that be] bowed down 雅威扶持一切跌倒的，扶起一切被壓下的。
\item \underline{Ain \<`>} The eyes all on wit You and You give like them their food in due time. 萬有的眼目都仰望你，你按時給他們食物。
\item \underline{Pe \<pp>} open Your hand, and satisfy of every living thing the desire 你張開你的手，使眾生都隨願飽足。
\item \underline{Tzaddi \<.s.s>} [is] righteous The LORD in all His ways, and holy in all his works. 雅威在祂一切所行的盡都公義，在祂一切所做的都有慈愛。
\item \underline{Koph \<q>} [is] near The LORD to all those who call To all that call upon Him in truth. 雅威臨近凡求告祂的，就是一切誠心求告祂的人。
\item \underline{Resh  \<r>} the desire of those who fear He will fulfill and their cry he also will hear and will save. 敬畏祂的人的心願，祂必成就；他們的呼求祂必垂聽，也必拯救他們。
\item \underline{Schin \<,s>} preserves The LORD, all those who love but the all wicked will he destroy. 雅威保護一切愛祂的人，卻要滅絕一切的惡人。
\item \underline{Tau \<t>} the praise of the LORD shall speak My mouth bless all flesh name his holy forever and ever. 我的口要述說雅威的讚美；惟願凡有血氣的都永永遠遠稱頌祂的聖名。
\end{itemize}




















